\section{Group determinants and the scalar Lehmer gap}\label{sec:determinants}

The determinant strategy seeks operator or geometric information that returns to the same scalar gap.
For the infinite cyclic group this return is exact.
More general groups, singular operators and rectangular matrices require their own comparison; their determinant constants cannot be silently substituted for the classical one \cite{Luc22}.

Let $S^1$ carry normalized arclength.
Fourier transformation identifies the group von Neumann algebra $\mathcal N(\Z)$ with the multiplication algebra $L^\infty(S^1)$ on $L^2(S^1)$.
A matrix $A\in M_r(\Z[t,t^{-1}])$ acts on $L^2(S^1)^r$ by multiplication by $A(e^{i\theta})$.

\begin{definition}
Such an operator is a weak isomorphism if it is injective and has dense image.
For a weak isomorphism, write $|A|=(A^*A)^{1/2}$ and let $\mu_{|A|}$ be its spectral measure for the circle trace and the unnormalized matrix trace $\operatorname{tr}_r(B)=\sum_{j=1}^rB_{jj}$.
Its Fuglede--Kadison determinant is
\[
 \det_{\mathcal N(\Z)}(A)
   =\exp\left(\int_{(0,\infty)}\log s\,d\mu_{|A|}(s)\right).
\]
\end{definition}
These are the weak-isomorphism and trace conventions of \citet{Luc22}.
In particular, the matrix trace is not divided by $r$.

\begin{proposition}\label{det:scalar-return}
For $A\in M_r(\Z[t,t^{-1}])$, the corresponding multiplication operator is a weak isomorphism exactly when $\det A$ is not identically zero.
In that case
\[
 \log\det_{\mathcal N(\Z)}(A)
  =\frac1{2\pi}\int_0^{2\pi}\log|\det A(e^{i\theta})|\,d\theta
  =m(\det A).
\]
For a Laurent polynomial, $m$ is defined by multiplying by a monomial to obtain an ordinary polynomial; the value is independent of that choice.
The infimum of the positive logarithmic determinants of these square weak isomorphisms, over all $r$, equals $\mathfrak L$.
\end{proposition}
\begin{proof}
A nonzero Laurent determinant vanishes at only finitely many points of the circle.
The matrix is consequently invertible almost everywhere, so both it and its adjoint have zero kernel.
Conversely an identically zero determinant gives a nontrivial measurable kernel.
At each invertible point, the unnormalized trace of $\log|A|$ is $\log|\det A|$, by the product of singular values.
Integrating proves the first equality; the scalar Mahler integral proves the second.
The integral is finite because a polynomial has only logarithmic singularities at its circle zeros.
Every determinant can be made an integer polynomial by a monomial shift, while every nonzero integer polynomial occurs as a $1$-by-$1$ example.
This proves the infimum claim.
\end{proof}

The normalization is essential.
Division by $r$ would produce $m(\det A)/r$, which changes the gap question when matrix size varies.
Proposition~\ref{det:scalar-return} also explains why the full scalar polynomial must be retained in an integer construction.

\begin{corollary}\label{det:full-cost}
If the full Laurent determinant of a square integer weak isomorphism is $c\,t^jP(t)Q(t)$, with $c\in\Z\setminus\{0\}$ and nonzero $P,Q\in\Z[t]$, then
\[
 \log\det_{\mathcal N(\Z)}(A)
   =\log|c|+m(P)+m(Q).
\]
Thus the complementary factor is an additional nonnegative expense.
\end{corollary}
\begin{proof}
Apply Proposition~\ref{det:scalar-return} and multiplicativity of Mahler measure.
A monomial has zero logarithmic measure.
\end{proof}

The literal signed two-sheet construction in Subsection~\ref{sec:structures} is a concrete application: its complete scalar is $p_Vp_{V^-}$, so the absolute companion cost $m(p_{V^-})$ is compulsory.
The repeated-cell theorem \ref{str:uniform-cost} pays this cost for every selected subset.
Proposition~\ref{str:changed-cost} shows that an actual change of bulk recurrence can make isolated cells neutral while complete couplings still cost more than $k\log2$.
These results inform both the completion and determinant strategies; the full constructions are stated once in Subsection~\ref{sec:structures}.

\begin{question}\label{det:completion-question}
Can an actual integral completion, change of basis or operator construction control every complementary scalar factor and all couplings at a useful uniform scale, while preserving the original small-measure input and its unnormalized scalar return?
\end{question}
An answer must specify the complete polynomial and the actual integral realization.
Neutral factors still affect rank and attachments even when they have zero Mahler measure.
Bounds on height or entry energy are separate from bounds on the full determinant.
The two established matrix-family obstructions leave other integral bases and completion mechanisms open; neither optimizes over the entire integer class.
